
\subsection{The general construction}
\label{ssec:relating-general}

We now give a general construction to show how two-step linearisation can
capture synchronisation linearisation. 
%
Fix a synchronisation object~|SyncObj|, a synchronisation specification object
|SyncSpec|, and synchronisation abstraction function~$\A$.  Let the operations
of the |SyncObj| be
\begin{scala}
  def op£\s i£(x: A£\s i£): B£\s i \qquad for $i = 1,\ldots,n$.£
\end{scala}

The two-step linearisation tester has operations $\op_i$ and
$\overline\op_i$ for $i = 1,\ldots,n$. 
%
The idea is that for a mapping 
\[
\seq{\op_{i_1}(x_1), \ldots, \op_{i_a}(x_a)}
  \mapsto \sync(x_1,\ldots,x_a)  \in \A,
\]
corresponding to a synchronisation of arity~$a$, is linearised in the
following order:
\begin{enumerate}
\item $\op_{i_1},\ldots,\op_{i_a}$, in order;

\item $\overline\op_{i_j}$, for some~$j$, which  calls
  |SyncSpec.sync|$(x_1,\ldots,x_a)$, and stores the results for other threads;

\item The remaining $\overline\op_{i_k}$, in some order, which  retrieve
  their results.
\end{enumerate}
%
No other operation occurs between the start of stage~1 until after stage~3
completes.  Thus operations for different synchronisations are not interleaved
(this isn't strictly necessary, but simplifies the subsequent proof). 

% stage~2; but other operations may occur during stage~3.

During stage~1, the two-step linearisability tester records the prefix of
$\seq{\op_{i_1}(x_1), \ldots, \op_{i_a}(x_a)}$ performed so far.  This is done
by defining a case class |Op|$_i$ corresponding to each such
operation~|op|$_i$, as at the top of Figure~\ref{fig:schema}.
%
%% \begin{scala}
%% trait Op
%% case class Op_i(t: ThreadID, x: A_i) £\quad for $i = 1,...,n$.£
%% type State = List[Op]
%% \end{scala}
%
The following function maps a value from such a class to the operation it
represents:
\begin{eqnarray*}
\opFor(\sm{Op}_i(t,x)) & = & op_i(x).
\end{eqnarray*}
%
We then define a function $\opsFor$ that maps a sequence of such $\sm{Op}_i$
values to the corresponding sequence of operations, by pointwise application
of $\opFor$.
%% The following function maps such a sequence to the corresponding sequence of
%% operations, for use with~$\A$.
%% \begin{eqnarray*}
%% \opsFor(state) &  = & \seq{\op_j(x_j) \midd \sm{Op}_j(t_j,x_j) \leftarrow state}.
%% \end{eqnarray*}
%

A schema (using pseudo-Scala) for the  construction is in
Figure~\ref{fig:schema}.\footnote{We write ``$\prefix$'' for sequence
  prefixing. ``{\scalastyle :+}'' is the Scala operator to append a value onto
  a list.}  The variable~|state| records the sequence of |op|$_i$ linearised
so far (more precisely, $\opsFor(\sm{state})$ is that sequence): this
satisfies an invariant that it is a prefix of an element of $\dom \A$.

%%%%%

\begin{figure}
\begin{scala}
trait Op
case class Op£$_i$£(t: ThreadID, x: A£$_i$£) //  for £$i = 1,...,n$.£
type State = List[Op]
type B = £$\bigcup\nolimits_{i = 1,\ldots,n} \sm B_i$£

class TwoStepLinSpec{
  private var state: State = List()
  private val returns = new Array[Option[B]](NumThreads)
  private var remaining = 0
  for(i <- 0 until NumThreads) returns(i) = None

  def op£\s i£(t: ThreadID, x: A£\s i£): Unit = { // for £$i = 1,...,n$£.
    require(remaining == 0 && £$\exists ops \in \dom \A \spot \opsFor(\sm{state}) \cat \seq{\op_i(x)} \prefix ops$£)
    state = state :+  Op£\s i£(t,x) 
  }

  def £$\overline\op_i$£(t: ThreadID): B£\s i£ = { // for £$i = 1,...,n$£.
    if(£$\opsFor(\sm{state}) \in \dom \A$£){
      val (y£\s 1£,...,y£$_{\ss a}$£) = SyncSpec.£$\A(\opsFor(\sm{state}))$£
      for(m <- 1 to a) returns(state(m).t) = Some(y£$_{\ss m}$£)
      remaining = a; state = List()
    }
    else (require(state == List() && returns(t).isInstanceOf[Some])
    val Some(y) = returns(t); returns(t) = None; remaining -= 1; y
  }
}
\end{scala}
\caption{A schema for the general construction of a two-step linearisability
  tester.} 
\label{fig:schema}
\end{figure}

%%%%%

For each thread~|t|,\, |returns(t)| stores the value to be returned by~|t| as
part of its most recent uncompleted synchronisation, if any.  The variable
|remaining| records how many values still need to be retrieved from~|returns|;
calls to $\op_i$ operations are blocked when this is non-zero, so as to
prevent interleaving of operations for different synchronisations.

The operation $\op_i$ appends an appropriate |Op|$_i$ object to~|state|.  The
precondition ensures that |state| still corresponds to a prefix of an element
of $\dom \A$.  

The operation $\overline\op_i$ proceeds as follows.  In the case that |state|
corresponds to an element of $\dom \A$, it calls the corresponding operation
on |SyncSpec|, stores each result in the corresponding entry of |returns|, and
resets |state|.  Otherwise, it requires that some other operation has called
the operation on |SyncSpec|, and that |state| holds its initial value.  In
either case, the operation retrieves and clears its result from the
appropriate entry in~|returns|.

%% It is then straightforward to check that operations of |TwoStepLinSpec| can be
%% linearised only as described earlier.  The proof that two-step linearisability
%% corresponds to synchronisation linearisability is then much as in the proof
%% for the binary case in Proposition~\ref{prop:two-step-lin}. 


%% We believe that these techniques can be adapted to all the other
%% synchronisation objects we have considered.

%% These techniques can be adapted to all other synchronisation objects of the
%% form we are considering.  We sketch a generic approach; we omit the details
%% because they are fiddly and uninteresting.  For each operation |op|$_i$ of the
%% synchronisation object, the two-step linearisation specification object
%% includes operations $\op_i$ and $\overline{\op}_i$.  For each maplet
%% $\seq{\op_{i_i}(x_1), \ldots, \op_{i_a}(x_a)} \mapsto \sync(x_1,\ldots,x_a)$
%% in the synchronisation abstraction function, the two-step linearisation
%% specification object performs operations as follows: (1)~$\op_{i_i}(x_1),
%% \ldots, \op_{i_a}(x_a)$, in that order; (2)~$\overline\op_{i_i}$, for
%% some~$i$, which calls \sync\ on the synchronisation linearisation object and
%% stores the results for other threads; (3)~the remaining $\overline\op_{i_j}$,
%% in some order, which retrieve their results.

%%%%%%%%%%

%\framebox{Work in progress}

The following definition identifies  operations of |TwoStepLinSpec| that
correspond to a single synchronisation.
%
\begin{definition}
A \emph{synchronisation subhistory} for event
\[
sync(x_1,\ldots,x_a)^{i_1,\ldots,i_a}\::(y_1,\ldots,y_a)
\]
is a sequence of events of |TwoStepLinSpec| of the form
\[
\seq{ \op^{i_1}_{j_1}(t_1, x_1)\::(),\; \ldots,\; 
  \op^{i_a}_{j_a}(t_a, x_a)\::(),\;
  \overline{\op}^{i_1'}_{j_1'}(t_1')\::y_1',\; \ldots,\; 
    \overline\op^{i_a'}_{j_a'}(t_a')\::y_a' }
\]
such that $(j_1',i_1',t_1',y_1'),\ldots,(j_a',i_a',t_a',y_a')$ is a
permutation of $(j_1,i_1,t_1,y_1),\linebreak[1] \ldots,\linebreak[1]
(j_a,i_a,t_a,y_a)$, and
\[
\A(\seq{\op_{j_1}(x_1), \ldots, \op_{j_a}(x_a)}) = sync(x_1,\ldots,x_a).
\]
\end{definition}
%
In other words, the threads call operations $\op_{j_1},\ldots,\op_{j_a}$, then
call the corresponding $\overline\op_j$ operations in some order, and receive
back the results $y_1,\ldots,y_a$; and the synchronisation abstraction
function associates these operation calls with~$sync(x_1,\ldots,x_a)$.  These
operations correspond to a single synchronisation consistent with
the synchronisation specification object~|SyncSpec| if
|SyncSpec.|$sync(x_1,\ldots,x_a)$ returns $(y_1,\ldots,y_a)$.  The
synchronisation can be considered to happen within~$\overline\op^{i_1'}_{j_1'}$.

%%%%%

%
\begin{definition}
Given a legal history $h_s$
%% \begin{eqnarray*}
%% h_s & = & \seq{ sync_1^{\vec{i_1}}(\vec{x_1})\::\vec{y_1}, \ldots,
%%   sync_k^{\vec{i_k}}(\vec{x_k})\::\vec{y_k} }
%% \end{eqnarray*}
of |SyncSpec|, a \emph{complete synchronisation history} for~$h_s$ is the
concenation of synchronisation subhistories for the events in $h_s$.
%% \begin{eqnarray*}
%% h_s & = & \seq{ sync_1^{\vec{i_1}}(\vec{x_1})\::\vec{y_1}, \ldots,
%%   sync_k^{\vec{i_k}}(\vec{x_k})\::\vec{y_k} }
%% \end{eqnarray*}
%% such that all the $\op_j$ events of one synchronisation subhistory are after
%% the first $\overline\op_j$ of  the previous.
%, and where $h_s$
% $\seq{sync_1(\vec{x_1})\::\vec{y_1}, \ldots, sync_k(\vec{x_k})\::\vec{y_k}}$
%is a legal history of |SyncSpec|.
\end{definition}
%
%% This represents several synchronisations, which may overlap to a certain
%% extent: $\op_i$ events of one synchronisation subhistory can occur once the
%% previous synchronisation has happened (in its first $\overline\op_i$ event);
%% but these $\op_i$ events may be interleaved with $\overline\op_i$ events of
%% previous synchronisation subhistories.

%% For example, assuming $\A(\seq{\op_1(x), \op_2(x')}) = \sync(x,x')$,
%% \[
%% \begin{align}
%% \seq{ \op_1^1(t_1,x_1)\::(),\; \op_2^2(t_2,x_2)\::(),\; 
%%   \overline\op_2^2(t_2)\::y_2,\;  \op_1^3(t_3, x_3)\::(), \\ 
%% \quad
%% \op_2^4(t_4,x_4)\::(),\; \overline\op_1^3(t_3)\::y_3,\;
%%  \overline\op_2^t(t_4)\::y_4,\; \overline\op_1^1(t_1)\::y_1 }
%% \end{align}
%% \]
%% is a complete synchronisation history for 
%% \[
%% \seq{\sync^{1,2}(x_1,x_2)\::(y_1,y_2),\;  \sync^{3,4}(x_3,x_4)\::(y_3,y_4)}.
%% \]
%% The final $\overline\op_1^1$ belongs to the first synchronisation, but is
%% delayed.

%%%%%

The following lemma follows directly from the construction of
|TwoStep|\-|LinSpec|.  
%
\begin{lemma}
\label{lem:TwoStepLinSpec}
Assume every thread alternates between calls of some~$\op_j$ and the
corresponding $\overline\op_j$ (passing in its own identity).
%
\begin{enumerate}
\item Suppose $h_s$ is a legal history of |TwoStepLinSpec| such that for
  every $\op_j^i$ event, there is a matching $\overline{\op}_j^i$ event; then
  $h_s$ is a complete synchronisation history (so corresponds to a legal
  history of~|SyncSpec|).

\item Every complete synchronisation history is a legal history of
  |TwoStepLinSpec|.
\end{enumerate}
\end{lemma}
%
Informally, |TwoStepLinSpec| allows $\op_j$ operations (adding them to
|state|); once the current sequence corresponds to a possible synchronisation,
then a call to some $\overline\op_i$ can trigger that, corresponding to the
relevant $sync$ event; and once all $\overline\op_j$
operations have retrieved their results, the $\op_j$ operations for the next
synchronisation can start.  Further, the calls to the operations on~|SyncSpec|
must correspond to a legal history of |SyncSpec|.

%% Note that this lemma allows for the case when some pending $\overline\op_i$
%% events have not yet happened. 

%% \framebox{Refine stuff above to give what we need below.}  I think we only
%% really need complete synchronisation histories.  Note that the linearisation
%% of a complete history of |SyncObj| is complete. 

Below, we denote a tuple of parameters $(x_1,\ldots,x_a)$ by a
vector $\vec{x}$; and likewise for return values and execution identities.
%
\begin{prop}
Each history~$h$ of~|SyncObj| is two-step-linearisable with respect to
|TwoStepLinSpec| if and only if it is synchronisation-linearisable with
respect to |SyncSpec|.  Hence |SyncObj| is two-step-linearisable with respect
to |TwoStepLinSpec| if and only if it is synchronisation-linearisable with
respect to |SyncSpec|.
\end{prop}
%
\begin{proof}
\textbf{($\implies$).}\quad
%
Let $h$ be a concurrent history of |SyncObj| that is two-step-linearisable
with respect to |Two|\-|Step|\-|LinSpec|.  By assumption, there is an
extension $h'$ of~$h$, and a legal history~$h_s$ of |TwoStepLinSpec| such that
$h_s$ is a two-step linearisation of $h'' = complete(h')$.

Necessarily, $h_s$ contains an $\overline{\op}_j^i$ event for each $\op_j^i$
event.  Hence, by Lemma~\ref{lem:TwoStepLinSpec}, $h_s$ is a complete
synchronisation history corresponding to a legal history $h_s'$ of~|SyncSpec|.
In effect, $h_s'$ replaces each synchronisation subhistory for event
$sync^{\vec{i}}(\vec{x})\::\vec{y}$ in~$h_s$ by the event
$sync^{\vec{i}}(\vec{x})\::\vec{y}$.
%% , where $\vec{i}$ is the sequence of
%% execution identities on the corresponding $\op$ events in~$h_s$.
%
This is illustrated by the example timeline below (for a single
synchronisation subhistory of arity~3), where $h''$ is represented by the
horizontal lines and the labels above; $h_s$~is represented by the
``$\cross$''s and the labels below; and $h_s'$ is represented by the
``$\bullet$'' and the label below.
%
\begin{center}
\begin{tikzpicture}[xscale = 0.8,yscale = 1]
\draw[|-|] (0,0) -- node[above] {$\op_1^1(x_1)\::y_1$} (7,0);
\crossAt(1,0){$\op_1^1(t_1,x_1):()$}
\crossAt(6,0){$\overline{\op_1^1}():y_1$}
%
\draw[|-|] (1.5,-1.4) -- node[above] {$\op_2^2(x_2)\::y_2$} (6,-1.4);
\crossAt(2.5,-1.4){$\op_2^2(t_2,x_2)\::()$};
\crossAt(5,-1.4){$\overline\op_2^2()\::y_2$};
%\draw(3,-1.5) \X; \draw(3,-1.9) node{$\sm{receive}^2()\::3$};
%
\draw[|-|] (2,-2.8) -- node[above] {$\op_3^3(x_3)\::y_3$} (7,-2.8);
\crossAt(3.0,-2.8){$\op_3^3(t_3,x_3)\::()$};
\crossAt(6.5,-2.8){$\overline\op_3^3()\::y_3$};
%
\draw(5,-3.8) node{$\bullet$}; 
\draw(5,-4.15) node{$\sm{sync}^{1,2,3}(x_1,x_2,x_3)\:: (y_1,y_2,y_3)$};
\end{tikzpicture}
\end{center}
%
%By Lemma~\ref{lem:TwoStepLinSpec}, $h_s'$ is a legal history of |SyncSpec|. 
It is possible to interleave $h''$ and~$h_s'$ by placing each
$sync^{\vec{i}}(\vec{x}) \:: \vec{y}$ event in the same place as the first
$\overline{\op}_j$ of its synchronisation subhistory in the interleaving
of~$h''$ and~$h_s$ (the $\overline{\op}_2$ event in the above diagram); by
construction, this is between each $\call.\sm{op}_j^{i_j}(x_j)$
and~$\return.\sm{op}_j^{i_j} \:: y_j$.
%
Hence $h_s'$ is a synchronisation linearisation of~$h''$; and so $h$ is
synchronisation-linearisable.

\textbf{($\Leftarrow$).}\quad
%
Let $h$ be a history of |SyncObj| that is synchronisation-linearisable with
respect to |SyncSpec|.  By assumption, there is an extension $h'$ of~$h$, and
a legal history~$h_s$ of |SyncSpec| such that $h_s$ is a synchronisation
linearisation of $h'' = complete(h')$.
%
Build a history~$h_s'$ of |TwoStepLinSpec| by replacing each  $sync$
event by a corresponding synchronisation subhistory.  More precisely, if the
event $sync^{i_1,\ldots,i_a}(x_1,\ldots,x_a) \:: (y_1,\ldots,y_a)$
corresponds to operation executions 
\(  \op_{j_1}^{i_1}(x_1)\::y_1, \ldots, \op_{j_a}^{i_a}(x_a)\::y_a \)
%
of |SyncObj|, 
with $\A(\seq{\op_{j_1}(x_1), \ldots, \op_{j_a}(x_a)}) = sync(x_1,\ldots,x_a)$,
then replace the $sync$ event by the consecutive events
%
\(  \op^{i_1}_{j_1}(t_1, x_1)\::(),\linebreak[1]\; \ldots,\; \op^{i_a}_{j_a}(t_a, x_a)\::(),\;
  \overline{\op}^{i_1}_{j_1}(t_1)\::y_1,\; \ldots,\;
  \overline\op^{i_a}_{l_a}(t_a)\::y_a  \).

This is a legal history of |TwoStepLinSpec|, by
Lemma~\ref{lem:TwoStepLinSpec}.  We can interleave $h''$ and~$h_s'$ by
inserting each synchronisation subhistory in the place of the corresponding
$sync$ event in the interleaving of $h''$ and~$h_s$.  By construction, each
$\sm{op}_{j_i}^{i_i}(t_i,x_i) \:: ()$ and $\overline{\sm{op}}_{j_i}^{i_i}()
\:: y_i$ are between $\call.\sm{op}_{j_i}^{i_i}(x_i)$
and~$\return.\sm{op}_{j_i}^{i_i} \:: y_i$.
%
Hence $h_s'$ is a two-step linearisation of $h''$; and so $h$ is
two-step-linearisable.
\end{proof}


Figure~\ref{fig:schema} uses pseudo-Scala.  For a concrete implementation, we
need a separate implementation of $\op_i$ and $\overline\op_i$ for each~$i$.
Then each calculation involving~$\A$ can be implemented directly, for example
using a case analysis.
